\section{An alternative construction from projection overlaps}\label{sec:projection}

The dimension-ten construction is complete.  Sections~\ref{sec:projection}
and \ref{sec:combinatorial} now produce another entangled member of
$\mathscr C$ in much larger finite dimensions.  The real construction has
an additional symmetry: the same trace-preserving rectangular PPT channel
acts on both halves of a maximally entangled state.  These are parallel
local actions on two subsystems, distinct from the self-composition of a
square channel in Theorem~\ref{thm:main-channel}.

The mechanism also differs from the dimension-ten range obstruction.
The input consists of real unit vectors and projective measurements.
The vectors produce a large sum of correlations, whereas
zero trace overlaps among the projections restrict the supports of the
product vectors in a hypothetical separable decomposition.  These two
bounds will conflict.  Trace Gram matrices of positive factors also occur
in completely positive semidefinite factorizations
\cite{PSVW2018}; here the zero pattern of such trace pairings supplies the
support restriction.

For self-adjoint projections \(E,F\), their trace overlap is always a
nonnegative real number:
\begin{equation}\label{eq:projection-overlap}
 \tr(EF)=\tr(FEF)=\tr\bigl((EF)^*EF\bigr)\geq0.
\end{equation}
In particular, it vanishes exactly when \(EF=0\).  This fact does not
require \(E\) and \(F\) to commute.

\begin{proposition}[Rectangular projection criterion]
\label{prop:projection-criterion}
Let \(Q,K,D\) be positive integers, let \(\mathcal Q\) be a set of
\(Q\) elements, and let \(u_q\in\mathbb R^K\), \(q\in\mathcal Q\), be
unit vectors.  Put \(c_{qr}=u_q^{\T}u_r\).  For each \(q\), let
\(\{P_i:i\in\mathcal I_q\}\) be a finite projective measurement on
\(\C^D\): its members are self-adjoint projections, they are mutually
orthogonal, and their sum is \(I_D\).  Zero projections are allowed.
Take the label sets \(\mathcal I_q\) disjoint, put
\(\mathcal I=\bigsqcup_{q\in\mathcal Q}\mathcal I_q\), and set
\(n=|\mathcal I|\).  Suppose that
\begin{enumerate}
 \item \(P_iP_j=P_jP_i\) whenever \(i\in\mathcal I_q\),
       \(j\in\mathcal I_r\), and \(c_{qr}\ne0\);
 \item there is a real number \(\kappa<Q/K\) such that every
       \(J\subseteq\mathcal I\) satisfying
       \(\tr(P_iP_j)>0\) for all distinct \(i,j\in J\) has
       \(|J|\leq\kappa\).
\end{enumerate}
After ordering \(\mathcal I\), define the complex-linear map
\(\Psi:\Mat n\to\Mat D\) on matrix units by
\begin{equation}\label{eq:projection-map}
 H_{ij}=c_{qr}P_iP_j,\qquad
 \Psi(E_{ij})=H_{ij}
 \quad(i\in\mathcal I_q,\ j\in\mathcal I_r).
\end{equation}
Then \(\Psi\) and its Hilbert--Schmidt adjoint
\(\Psi^\dagger:\Mat D\to\Mat n\) are PPT.  Moreover,
\begin{equation}\label{eq:projection-composite}
 C:=J(\Psi^\dagger\circ\Psi)
 \quad\text{is nonseparable over }\C^n\otimes\C^n,
 \qquad \tr C=Q^2D.
\end{equation}
\end{proposition}

\begin{proof}
We first prove that $\Psi$ is PPT.  To establish nonseparability, let
$B$ be the compression of $C$ to the span of the vectors
$e_i\otimes e_i$, expressed in that basis, and let
$\mathbf1=(1,\ldots,1)^\T$.  We will compare two bounds on the sum
of its entries.  The vector Gram matrix gives
\[
 \mathbf1^*B\mathbf1\geq\frac QK\tr B,
\]
whereas the zero tensor-basis diagonal entries would force a separable
decomposition of $C$ to satisfy
$\mathbf1^*B\mathbf1\leq\kappa\tr B$.
We will also have $\tr B=QD>0$, so the strict inequality $\kappa<Q/K$
will make these bounds incompatible.

\emph{The block Gram matrix and its partial transposes.}
Regard each \(u_q\) also as a vector in \(\C^K\).  For arbitrary
\(w_i\in\C^D\), self-adjointness of the projections gives
\begin{equation}\label{eq:projection-block-gram}
 \sum_{i,j\in\mathcal I}\langle w_i,H_{ij}w_j\rangle
 =
 \left\|
   \sum_{q\in\mathcal Q}\sum_{i\in\mathcal I_q}u_q\otimes P_iw_i
 \right\|^2\geq0.
\end{equation}
Thus the block matrix
\(H=[H_{ij}]_{i,j\in\mathcal I}=J(\Psi)\) is positive semidefinite.
Lemma~\ref{lem:choi-separable} implies that \(\Psi\) is CP.

Since \(c_{qr}=c_{rq}\) is real, the first hypothesis gives
\begin{equation}\label{eq:projection-block-symmetry}
 H_{ij}=H_{ji}=H_{ij}^*.
\end{equation}
Indeed, if \(c_{qr}\ne0\), the two projections commute; if
\(c_{qr}=0\), both blocks vanish and no commutation is needed.
Let \(\Gamma_1=T_n\otimes\mathrm{id}_D\) and
\(\Gamma_2=\mathrm{id}_n\otimes T_D\) denote the partial transposes
on the input and output Choi factors.  Block symmetry gives
\begin{equation}\label{eq:projection-partial-transposes}
 H^{\Gamma_1}=H,\qquad
 H^{\Gamma_2}=H^{\T}\succeq0.
\end{equation}
For the second identity, full transpose is the product of the two
partial transposes.  Full transpose preserves complex positive
semidefiniteness: from \(H=R^*R\) we get
\(H^{\T}=(\overline R)^*\overline R\).  The Choi identities in
Lemma~\ref{lem:choi-separable} now show that \(T_D\circ\Psi\) is CP.
Hence \(\Psi\) is PPT.

Lemma~\ref{lem:ppt-stability} gives complete positivity and the PPT
property for \(\Psi^\dagger\).  We will also use a more specific
identity.  Testing the defining adjoint relation on matrix units gives,
for every \(A\in\Mat D\),
\begin{equation}\label{eq:projection-adjoint}
 [\Psi^\dagger(A)]_{ij}=\tr(H_{ij}^*A).
\end{equation}
Equation~\eqref{eq:projection-block-symmetry} therefore yields
\begin{equation}\label{eq:projection-adjoint-transpose}
 T_n\circ\Psi^\dagger=\Psi^\dagger
\end{equation}
on all complex matrices, including non-Hermitian ones.

\emph{The tensor-basis diagonal and a compression.}
The composition \(\Psi^\dagger\circ\Psi\) is CP, so \(C\succeq0\).
Index its rows and columns by \(\mathcal I\times\mathcal I\).
Equation~\eqref{eq:projection-adjoint} gives the fully indexed formula
\begin{equation}\label{eq:projection-composite-entries}
 C_{(i,a),(j,b)}=\tr(H_{ab}^*H_{ij}).
\end{equation}
Since \(c_{qq}=1\) and \(P_i^2=P_i\), we have \(H_{ii}=P_i\).
The diagonal entries with indices \((i,j)\), and the entries of the
compression \(B\), are consequently
\begin{align}
 C_{(i,j),(i,j)}&=\tr(P_iP_j),
       \label{eq:projection-crossed-diagonal}\\
 B_{ij}:=C_{(i,i),(j,j)}
   &=\tr(H_{ij}^*H_{ij})
     =c_{qr}^{\,2}\tr(P_iP_j)
       \quad(i\in\mathcal I_q,\ j\in\mathcal I_r).
       \label{eq:projection-compression}
\end{align}
The last equality follows from
\(\tr(P_jP_iP_iP_j)=\tr(P_iP_j)\), by idempotence and cyclicity of
trace.  It remains valid without commutation and when \(c_{qr}=0\).

For each pair \(q,r\), completeness of the two measurements gives
\[
 \sum_{i\in\mathcal I_q}\sum_{j\in\mathcal I_r}\tr(P_iP_j)
 =\tr\!\left[
    \left(\sum_{i\in\mathcal I_q}P_i\right)
    \left(\sum_{j\in\mathcal I_r}P_j\right)
   \right]=D.
\]
Summing \eqref{eq:projection-crossed-diagonal} proves
\(\tr C=Q^2D\).  The same identity and
\eqref{eq:projection-compression} give
\begin{equation}\label{eq:projection-gram-bound}
 \tr B=QD,\qquad
 \mathbf1^*B\mathbf1
   =D\sum_{q,r\in\mathcal Q}c_{qr}^{\,2}
   \geq\frac{DQ^2}{K}.
\end{equation}
To see the inequality, set
\(M=\sum_{q\in\mathcal Q}u_qu_q^{\T}\), a positive semidefinite
\(K\)-by-\(K\) real matrix.  Then
\(\tr M=Q\) and \(\tr(M^2)=\sum_{q,r}c_{qr}^2\).
Cauchy--Schwarz applied to its \(K\) nonnegative eigenvalues yields
\(\tr(M^2)\geq(\tr M)^2/K=Q^2/K\).  Equivalently, the vector Gram
matrix has rank at most \(K\).

\emph{The support bound forced by separability.}
Suppose that \(C\) is separable over \(\C^n\otimes\C^n\).
Spectrally decomposing its local positive factors and absorbing positive
weights into the vectors gives a finite decomposition
\[
 C=\sum_\ell
   (x^{(\ell)}\otimes y^{(\ell)})
   (x^{(\ell)}\otimes y^{(\ell)})^*,
 \qquad x^{(\ell)},y^{(\ell)}\in\C^n.
\]
Whenever \(\tr(P_iP_j)=0\), the corresponding tensor-basis diagonal entry is
\[
 0=C_{(i,j),(i,j)}
   =\sum_\ell |x_i^{(\ell)}y_j^{(\ell)}|^2.
\]
Each summand is nonnegative, so \(x_i^{(\ell)}y_j^{(\ell)}=0\) for
every \(\ell\).  For one summand put \(z_i=x_i y_i\).
If distinct \(i,j\) lie in \(\supp z\), then \(x_i y_j\ne0\).
The preceding implication and \eqref{eq:projection-overlap} force
\(\tr(P_iP_j)>0\).  The second hypothesis therefore gives
\(|\supp z|\leq\kappa\).

Compressing the separable decomposition to the diagonal tensor
subspace gives \(B=\sum_\ell z^{(\ell)}(z^{(\ell)})^*\).
Complex Cauchy--Schwarz on each support now yields
\begin{equation}
 \begin{aligned}
 \mathbf1^*B\mathbf1
  &=\sum_\ell\left|\sum_i z_i^{(\ell)}\right|^2
    \leq \kappa\sum_\ell\|z^{(\ell)}\|^2
     =\kappa\tr B\\
  &<\frac{Q}{K}\tr B
     =\frac{DQ^2}{K}.
 \end{aligned}
 \label{eq:projection-separable-bound}
\end{equation}
The inequality is also valid for a zero \(z^{(\ell)}\), whose support
is empty.  The strict step uses \(\kappa<Q/K\) and
\(\tr B=QD>0\).  This contradicts
\eqref{eq:projection-gram-bound} and proves nonseparability.
\end{proof}

The complex criterion has a state interpretation
without the reality assumption.  By
\eqref{eq:projection-adjoint-transpose}, Lemma~\ref{lem:choi-transfer}
gives
\begin{equation}\label{eq:projection-complex-transfer}
 C=
 \bigl((\Psi^\dagger\circ T_D)\otimes\Psi^\dagger\bigr)
       (\Omega_D\Omega_D^*).
\end{equation}
Both local maps in this formula satisfy the input-transpose CP
conditions, as the transfer lemma asserts for the PPT pair
\(\Psi,\Psi^\dagger\).  Thus \(C/(Q^2D)\in\mathscr C\) also in the
complex case, and Theorem~\ref{thm:zero-key} and
Proposition~\ref{prop:operational-extensions} apply to it.  The two local
maps in \eqref{eq:projection-complex-transfer} may differ; the reality
assumption is what gives the same channel in
\eqref{eq:projection-real-state}.

\begin{corollary}[The same local channel for real projections]
\label{cor:real-channel}
Under the hypotheses of Proposition~\ref{prop:projection-criterion},
suppose that every \(P_i\) is real symmetric in the fixed basis of
\(\C^D\), and put \(\Lambda=\Psi^\dagger:\Mat D\to\Mat n\).
Then, on all complex inputs,
\begin{equation}\label{eq:projection-real-invariances}
 \Lambda\circ T_D=\Lambda=T_n\circ\Lambda,
 \qquad \tr\Lambda(A)=Q\,\tr A.
\end{equation}
In particular, \(\Lambda/Q\) is a trace-preserving PPT channel.
For \(\Phi_D=\Omega_D/\sqrt D\), the normalized maximally entangled
vector on \(\C^D\otimes\C^D\), one has
\begin{equation}\label{eq:projection-real-state}
 \frac{C}{Q^2D}
 =
 \left(\frac{\Lambda}{Q}\otimes\frac{\Lambda}{Q}\right)
       (\Phi_D\Phi_D^*).
\end{equation}
This is an entangled density operator on \(\C^n\otimes\C^n\) in
\(\mathscr C\).
\end{corollary}

\begin{proof}
Reality and \eqref{eq:projection-block-symmetry} give the additional
identity \(H_{ij}^{\T}=H_{ij}\).  Equation~\eqref{eq:projection-adjoint}
therefore becomes
\[
 [\Lambda(A)]_{ij}=\tr(H_{ij}A),\qquad
 [\Lambda(T_D(A))]_{ij}=\tr(H_{ij}^{\T}A)=\tr(H_{ij}A).
\]
This proves the first transpose invariance; the second is
\eqref{eq:projection-adjoint-transpose}.  Also,
\[
 \tr\Lambda(A)=\sum_i\tr(H_{ii}A)
 =\sum_{q\in\mathcal Q}\sum_{i\in\mathcal I_q}\tr(P_iA)
 =Q\,\tr A.
\]
Since \(\Lambda\) is CP and PPT by the proposition, \(\Lambda/Q\)
is a trace-preserving PPT channel with input size \(D\) and output
size \(n\).

Apply Lemma~\ref{lem:choi-transfer} to the pair \(\Psi,\Lambda\).
Its first local factor is
\[
 \Psi^\sharp=T_n\circ\Lambda\circ T_D=\Lambda
\]
by the two invariances.  Thus
\begin{equation}\label{eq:projection-real-identification}
 G:=(\Lambda\otimes\Lambda)(\Omega_D\Omega_D^*)=C.
\end{equation}
Dividing this identity by \(Q^2D\) gives
\eqref{eq:projection-real-state}.  Proposition~\ref{prop:projection-criterion}
gives \(\tr C=Q^2D>0\) and complex nonseparability.  Finally,
\(\Lambda/Q\) and \((\Lambda/Q)\circ T_D\) are CP, so
\eqref{eq:projection-real-state} is a representation of the form in
Definition~\ref{def:represented-class}.  This proves the assertion.
\end{proof}

The two channel uses in \eqref{eq:projection-real-state} act on separate
subsystems.  We next embed the rectangular maps into one matrix algebra
to obtain a composition of square PPT maps.

\begin{corollary}[Square embedding]\label{cor:square-embedding}
Under the hypotheses of Proposition~\ref{prop:projection-criterion},
suppose in addition that \(n\geq D\), and let
\(U:\C^D\to\C^n\) be the real coordinate isometry onto the first
\(D\) coordinates.  Define
\begin{equation}\label{eq:projection-square-pair}
 \Xi_1=\Ad_U\circ\Psi,\qquad
 \Xi_2=\Psi^\dagger\circ\Ad_{U^*}
       =\Xi_1^\dagger .
\end{equation}
Then \(\Xi_1,\Xi_2:\Mat n\to\Mat n\) are PPT and
\[
 \Xi_2\circ\Xi_1=\Psi^\dagger\circ\Psi.
\]
Their composition is not entanglement breaking.
\end{corollary}

\begin{proof}
Embedding and compression by \(U\) are CP.  The PPT composition
stability in Lemma~\ref{lem:ppt-stability} therefore applies to both
maps.  Directly, the reality of \(U\) also gives
\(T_n(UAU^*)=U T_D(A)U^*\), while
\eqref{eq:projection-adjoint-transpose} fixes the output transpose of
the second map.  The adjoint identity follows by reversing the two
factors, and \(U^*U=I_D\) gives the composite identity.  Its Choi
matrix is the nonseparable \(C\) of
Proposition~\ref{prop:projection-criterion}; the characterization in
Lemma~\ref{lem:choi-separable} proves the last assertion.
\end{proof}

The dimension comparison in this corollary is used only to choose
the isometry \(U\).  The rectangular criterion itself has no such
comparison.  The displayed square maps need not preserve trace or be
unital.
